\scanpage{21}

\[
\mathrm{Spec}\,I_k=\mathrm{Spec}\,k^I=\coprod_I\underbrace{\mathrm{Spec}(k)}_{S}
\qquad\text{\unsure{envisagé|envoyé} dans \unsure{Sch|$S$} de façon évidente.}
\]

b) Produits fibrés dans le cas où l'un des deux morphismes est un monomorphisme. Alors le diagramme correspondant sur les spectres est cartésien

\diagram{\begin{tikzcd}[ampersand replacement=\&,row sep=small]
\mathscr X\arrow[r,dashed,leftrightarrow]\arrow[d]\&\mathscr X'\arrow[d,dotted]\\
\mathscr Y\arrow[r,leftrightarrow]\&\mathscr Y'
\end{tikzcd}}

\diagram{\begin{tikzcd}[ampersand replacement=\&,row sep=small]
X\arrow[r]\arrow[d]\&X'\arrow[d]\\
Y\arrow[r]\&Y'
\end{tikzcd}}

De plus, si \uline{$Y'\to Y$} est une imm. fermée (resp. ouverte) -- plus généralement si cela correspond à $A\to AS^{-1}$, donc à un homéomorphisme dans, alors \deleted{le diagramme} de $=$ pour $\mathscr X'\unsure{\leftarrow|-|\to}\mathscr X$, donc la top. de $X'$ est la top. induite.

\uline{Cas particulier} Intersection de deux sous-objets ---

\diagram{\begin{tikzcd}[ampersand replacement=\&,row sep=small]
c)\quad A\arrow[r,"u"]\& B\\
Y\&X\arrow[l,"\varphi"]
\end{tikzcd}}

formule pour $\overline{\varphi(Z)}$, $Z=V(J)$

Condition pour que $\varphi$ soit dominant :

\[
\mathrm{Ker}\,\varphi\subset\mathrm{Nil}\,A
\]

% archive sheet number: 20
